
\Problem{Sous desorbitats}

\Statement

La gestió nefasta d'un club esportiu ha provocat que alguns dels seus
esportistes professionals tinguin sous desorbitats. Implementeu la
funció
\begin{lstlisting}
void split(const BD& data, const Fitxa& f, vector<string>& up, vector<string>& down);
\end{lstlisting}
que, a partir d'una base de dades @data@ que conté les fitxes dels
esportistes ordenades decreixenment per nivell d'absurditat, i donada
la fitxa @f@ d'un d'aquests esportistes, retorna dos vectors de textos
@up@ i @down@, on @up@ conté els noms dels esportistes amb fitxes mes
absurdes que @f@, i @down@ conté els noms dels esportistes amb fitxes
menys absurdes que @f@, ambdós en el mateix ordre en què apareixen a
@data@.  \bigskip

Cada fitxa consta d'un identificador únic (un enter), el nom de
l'esportista (un text), i una llista no buida de rebuts de nòmina (un
vector d'enters).

\begin{lstlisting}
struct Fitxa {
    int id;
    string nom;
    vector<int> nomines;
}
\end{lstlisting}


Es considera que una fitxa @f1@ és \emph{mes absurda} que una fitxa
@f2@ si, o bé el sou total de @f1@ és mes gran que el sou total de
@f2@, o bé els sous totals són iguals però @f1@ conté menys rebuts de
nòmina que @f2@, o bé coincideixen en aquestes dues coses però les
fitxes no tenen exactament els mateixos rebuts de nòmina i, al primer
rebut en què les quantitats difereixen, la quantitat de @f1@ és mes
gran que la de @f2@. En cas que les fitxes coincideixin en tots els
rebuts de nòmina, es considera que @f1@ és mes absurda que @f2@ si
l'identificador de @f1@ és mes petit que l'identificador de @f2@.
\bigskip

La base de dades és un vector de fitxes ordenat decreixenment per
nivell d'absurditat:

\begin{lstlisting}
using BD = vector<Fitxa>;
\end{lstlisting}


Per ajudar-vos a resoldre el problema, descarregueu el fitxer \texttt{code.cc}
disponible a la pàgina de l'enunciat. Hi trobareu el programa principal \textit{main()}
que s'encarrega de tota la gerstió de l'entrada/sortida i de la crida a \textit{split()}.
Per fer proves, podeu fer servir
l'exemple que us proposem. La sortida conté el número total de fitxes,
el número de fitxes mes absurdes que la fitxa donada, el número de
fitxes menys absurdes que la fitxa donada, seguit dels noms a @up@ i
@down@.

\medskip

A més, també us oferim un tros de codi que implementa una cerca dicotòmica en un vector d'enters ordenat creixentment. Si ho trobeu necessari, adapteu-lo al vostre problema.

\clearpage
\SampleOneCol
